% ===========================================================================
\section{La gramática original, comentada}
% ===========================================================================

Antes de tocar nada, un recorrido por la gramática de la cátedra explicando qué modela cada regla
y \textbf{por qué está escrita así}. Varias decisiones de forma que parecen arbitrarias son en
realidad restricciones del lenguaje impuestas por la gramática, y perderlas al transformar sería
cambiar \textsc{MiniJava}.

\subsection{Estructura del programa}

\begin{producciones}
\regla{Inicial}{\nt{ListaClases} \tl{eof}}
\regla{ListaClases}{\nt{Clase} \nt{ListaClases} \alt{} \nt{Interfaz} \nt{ListaClases} \alt{} \eps}
\end{producciones}

El \tl{eof} no es decoración: cumple la función del \tl{\$} de la teoría y obliga al analizador a
consumir \emph{todo} el archivo. Sin él, un programa con símbolos sobrantes después de la última
clase se aceptaría, porque \nt{ListaClases} puede derivar a vacío y el análisis terminaría dejando
tokens sin leer. En el léxico de la Etapa 1, \tl{eof} es un token con lexema \texttt{\$}.

\nt{ListaClases} usa \textbf{recursión a derecha}, que es la forma natural de expresar «cero o
más» en una gramática destinada al análisis descendente. La variante con recursión a izquierda
generaría el mismo lenguaje y sería inservible.

\begin{producciones}
\regla{Clase}{\tl{class} \tl{idClase} \nt{GenericidadOpcional} \nt{HerenciaOpcional} \tl{\{} \nt{ListaMiembros} \tl{\}}}
\regla{Interfaz}{\tl{interface} \tl{idClase} \nt{GenericidadOpcional} \nt{ExtensionOpcional} \tl{\{} \nt{ListaMetodosInterfaz} \tl{\}}}
\regla{GenericidadOpcional}{\tl{<} \tl{idGen} \tl{>} \alt{} \eps}
\regla{HerenciaOpcional}{\tl{extends} \nt{TipoReferencia} \alt{} \tl{implements} \nt{TipoReferencia} \alt{} \eps}
\regla{ExtensionOpcional}{\tl{extends} \nt{TipoReferencia} \alt{} \eps}
\end{producciones}

Clase e interfaz son casi simétricas, con tres diferencias deliberadas. Primero, la interfaz sólo
admite \tl{extends} y la clase admite \tl{extends} o \tl{implements}: por eso hay dos no terminales
distintos en vez de uno compartido, y la restricción queda impuesta \emph{por la gramática} en
lugar de delegarse al análisis semántico. Segundo, la interfaz sólo contiene métodos. Tercero,
\nt{HerenciaOpcional} tiene tres alternativas y no una combinación, lo que impide
\texttt{extends A implements B} --- válido en Java, no en \textsc{MiniJava}.

\nt{GenericidadOpcional} admite un \textbf{único} \tl{idGen}, no una lista: \textsc{MiniJava}
limita a un solo parámetro de tipo, y el logro \emph{Genericidad Avanzada E2} confirma que esa
limitación se mantiene.

\subsection{Miembros}

\begin{producciones}
\regla{Miembro}{\nt{Atributo} \alt{} \nt{Metodo} \alt{} \nt{Constructor}}
\regla{Atributo}{\nt{Tipo} \tl{idMetVar} \tl{;}}
\regla{Metodo}{\nt{ModificadorOpcional} \nt{TipoMetodo} \tl{idMetVar} \nt{ArgsFormales} \nt{Bloque}}
\regla{MetodoInterfaz}{\nt{TipoMetodo} \tl{idMetVar} \nt{ArgsFormales} \tl{;}}
\regla{Constructor}{\tl{public} \tl{idClase} \nt{ArgsFormales} \nt{Bloque}}
\regla{ModificadorOpcional}{\tl{static} \alt{} \eps}
\regla{TipoMetodo}{\nt{Tipo} \alt{} \tl{void}}
\end{producciones}

El \tl{public} del constructor \textbf{no es un modificador de visibilidad}: es lo que
sintácticamente lo distingue de todo lo demás. Por eso ni el atributo ni el método lo llevan. Este
detalle importa: el logro \emph{Visibilidad Mejorada E2} le saca esa función al \tl{public} y con
eso crea un conflicto nuevo (sección~\ref{sec:logro-visibilidad}).

\nt{Tipo} y \nt{TipoMetodo} están separados sólo para que \tl{void} sea un tipo de retorno válido
y no un tipo de variable válido: así \texttt{void x;} es imposible sin ningún chequeo adicional.

\nt{MetodoInterfaz} termina en \tl{;} y no tiene cuerpo. La diferencia con \nt{Metodo} está en el
último símbolo, y como se usan desde no terminales distintos nunca compiten.

\subsection{Tipos}

\begin{producciones}
\regla{Tipo}{\nt{TipoBase} \nt{DimensionesOpcionales}}
\regla{TipoBase}{\nt{TipoPrimitivo} \alt{} \nt{TipoReferencia} \alt{} \tl{idGen}}
\regla{DimensionesOpcionales}{\tl{[} \tl{]} \nt{DimensionesOpcionales} \alt{} \eps}
\regla{TipoReferencia}{\tl{idClase} \nt{TipoGenericoOpcional}}
\regla{TipoPrimitivo}{\tl{boolean} \alt{} \tl{char} \alt{} \tl{int}}
\regla{TipoGenericoOpcional}{\tl{<} \nt{InstanciadoOParametrico} \tl{>} \alt{} \eps}
\regla{InstanciadoOParametrico}{\tl{idGen} \alt{} \tl{idClase}}
\end{producciones}

Las tres familias de \nt{TipoBase} se distinguen con un solo token porque \textbf{el léxico ya
separó} \tl{idClase} de \tl{idGen} --- un \tl{idGen} es una única letra mayúscula. Es un buen
ejemplo de una decisión léxica que le simplifica la vida al analizador sintáctico.

\nt{InstanciadoOParametrico} nombra los dos usos de los ángulos: \emph{instanciado}
(\texttt{Lista<Persona>}, se fija un tipo concreto) y \emph{paramétrico} (\texttt{Lista<T>} dentro
de una clase genérica, se propaga el parámetro). No admite anidamiento ni notación diamante: las
dos cosas las agrega un logro.

\begin{decision}{transcripción de \texttt{[]}}
El documento de la cátedra escribe \tl{[]} como si fuera un símbolo, pero léxicamente son
\textbf{dos} tokens: \emph{Aspectos Léxicos de MiniJava} no define ningún token
\texttt{[]}. En este informe se escribe \tl{[} \tl{]}, que es lo que el analizador tiene que
reconocer. Es la única desviación deliberada de la literalidad del documento original.
\end{decision}

\subsection{Sentencias}

\begin{producciones}
\regla{Bloque}{\tl{\{} \nt{ListaSentencias} \tl{\}}}
\regla{Sentencia}{\tl{;} \alt{} \nt{Asignacion} \tl{;} \alt{} \nt{Llamada} \tl{;} \alt{} \nt{VarLocal} \tl{;} \alt{} \nt{Return} \tl{;}}
\regla{Sentencia}{\nt{If} \alt{} \nt{While} \alt{} \nt{Bloque}}
\regla{Asignacion}{\nt{Expresion}}
\regla{Llamada}{\nt{Expresion}}
\regla{VarLocal}{\tl{var} \tl{idMetVar} \tl{=} \nt{ExpresionCompuesta}}
\regla{Return}{\tl{return} \nt{ExpresionOpcional}}
\regla{If}{\tl{if} \tl{(} \nt{Expresion} \tl{)} \nt{Sentencia}}
\regla{If}{\tl{if} \tl{(} \nt{Expresion} \tl{)} \nt{Sentencia} \tl{else} \nt{Sentencia}}
\regla{While}{\tl{while} \tl{(} \nt{Expresion} \tl{)} \nt{Sentencia}}
\end{producciones}

Las cuatro primeras alternativas de \nt{Sentencia} llevan el \tl{;} explícito para que quede
visible en la regla; las tres últimas no lo llevan porque terminan en \tl{\}} o en otra sentencia.
La alternativa \tl{;} sola es la sentencia vacía.

\nt{VarLocal} declara \textbf{sólo con} \tl{var} y \textbf{siempre} con inicialización: en la
\textsc{MiniJava} base no existe \texttt{int x;}. Además usa \nt{ExpresionCompuesta} y no
\nt{Expresion} a la derecha del \tl{=}, lo que impide \texttt{var x = y = 3;}. Es una regla más
restrictiva a propósito.

\subsection{Expresiones}

Conviene leerlas como una pila de capas, cada una llamando a la de abajo. Ese encadenamiento es lo
que fija qué se agrupa con qué.

\begin{producciones}
\regla{Expresion}{\nt{ExpresionCompuesta} \nt{OperadorAsignacion} \nt{ExpresionCompuesta}}
\regla{Expresion}{\nt{ExpresionCompuesta}}
\regla{OperadorAsignacion}{\tl{=}}
\regla{ExpresionCompuesta}{\nt{ExpresionCompuesta} \nt{OperadorBinario} \nt{ExpresionBasica}}
\regla{ExpresionCompuesta}{\nt{ExpresionBasica}}
\regla{OperadorBinario}{\tl{||} \alt{} \tl{\&\&} \alt{} \tl{==} \alt{} \tl{!=} \alt{} \tl{<} \alt{} \tl{>} \alt{} \tl{<=} \alt{} \tl{>=} \alt{} \tl{+} \alt{} \tl{-} \alt{} \tl{*} \alt{} \tl{/} \alt{} \tl{\%}}
\regla{ExpresionBasica}{\nt{OperadorUnario} \nt{Operando} \alt{} \nt{Operando}}
\regla{OperadorUnario}{\tl{+} \alt{} \tl{-} \alt{} \tl{!}}
\regla{Operando}{\nt{Primitivo} \alt{} \nt{Referencia}}
\regla{Primitivo}{\tl{true} \alt{} \tl{false} \alt{} \tl{intLiteral} \alt{} \tl{charLiteral} \alt{} \tl{null}}
\end{producciones}

Dos observaciones que condicionan todo lo que sigue.

\textbf{La asignación no es asociativa.} El lado derecho del \tl{=} es una \nt{ExpresionCompuesta},
no una \nt{Expresion}. En consecuencia \texttt{a = b = c} \textbf{no es un programa
MiniJava válido}. Es fácil pasarlo por alto y «mejorarlo» al transformar; hacerlo
agrandaría el lenguaje.

\textbf{No hay niveles de precedencia entre operadores binarios.} Los trece operadores están en una
única alternancia plana, sin la jerarquía habitual de un no terminal por nivel. Por lo tanto
\texttt{a + b * c} y \texttt{a * b + c} se derivan con la misma forma: la precedencia y la
asociatividad no están en la gramática de \textsc{MiniJava}, y darles lugar sería cambiar el
lenguaje. Como corolario, eliminar la recursión a izquierda de \nt{ExpresionCompuesta} no pierde
nada, porque no había precedencia que perder.

\subsection{Referencias y primarios}

\begin{producciones}
\regla{Referencia}{\nt{Primario} \alt{} \nt{Referencia} \nt{VarEncadenada} \alt{} \nt{Referencia} \nt{MetodoEncadenado} \alt{} \nt{Referencia} \nt{AccesoArreglo}}
\regla{Primario}{\tl{this} \alt{} \tl{stringLiteral} \alt{} \nt{AccesoVar} \alt{} \nt{CrearArreglo}}
\regla{Primario}{\nt{LlamadaConstructor} \alt{} \nt{LlamadaMetodo} \alt{} \nt{LlamadaMetodoEstatico} \alt{} \nt{ExpresionParentizada}}
\regla{AccesoVar}{\tl{idMetVar}}
\regla{LlamadaMetodo}{\tl{idMetVar} \nt{ArgsActuales}}
\regla{LlamadaMetodoEstatico}{\tl{idClase} \tl{.} \tl{idMetVar} \nt{ArgsActuales}}
\regla{CrearArreglo}{\tl{new} \nt{TipoBase} \nt{DimensionesConTamanio}}
\regla{LlamadaConstructor}{\tl{new} \nt{TipoReferencia} \nt{ArgsActuales}}
\regla{ExpresionParentizada}{\tl{(} \nt{Expresion} \tl{)}}
\regla{VarEncadenada}{\tl{.} \tl{idMetVar}}
\regla{MetodoEncadenado}{\tl{.} \tl{idMetVar} \nt{ArgsActuales}}
\regla{AccesoArreglo}{\tl{[} \nt{Expresion} \tl{]}}
\end{producciones}

\nt{Referencia} modela una cadena como \texttt{this.lista.get(3).nombre[0]}: se arranca de un
\nt{Primario} y se le cuelgan eslabones. La recursión a izquierda expresa que la cadena se asocia
hacia la izquierda, que es lo correcto: \texttt{a.b.c} es \texttt{(a.b).c}.

\tl{stringLiteral} está en \nt{Primario} y no en \nt{Primitivo} \textbf{a propósito}: de un
\emph{String} se pueden colgar eslabones (\texttt{"hola".length()}) y de un \tl{intLiteral} no.
Por el mismo motivo está \nt{ExpresionParentizada}, que habilita \texttt{(a + b).metodo()}.

Ningún \nt{Primario} deriva un \tl{idClase} suelto: no existe la referencia a una clase por sí
sola, siempre tiene que ser \texttt{C.metodo(\dots)}. Ese hecho se usa más adelante.

\subsection{Construcciones que faltan}

Tres ausencias, que no son errores sino el espacio que dejan los logros:

\begin{itemize}
  \item \tl{++} y \tl{-{}-} son tokens del léxico --- \emph{Aspectos Léxicos} \S2.4.2 los lista ---
        pero \textbf{no aparecen en ninguna regla}. \nt{OperadorUnario} es sólo
        \tl{+} \alt{} \tl{-} \alt{} \tl{!}. Son tokens sin lugar hasta el logro
        \emph{Operadores Posifijos E2}.
  \item \tl{float} y \tl{floatLiteral} tampoco tienen lugar, y el léxico de la Etapa 1 los emite
        por un logro de aquella etapa. Sin extender la gramática, ese logro queda inutilizable.
  \item \tl{for}, \tl{private}, \tl{?} y \tl{->} no están \textbf{ni en la gramática ni en el
        léxico}: los agregan cuatro de los logros de esta etapa, y por lo tanto exigen tocar
        también el analizador léxico.
\end{itemize}
